
\chapter{Formulação da cinemática direta}
\label{ch:dir_kin}

\hspace*{\parindent}A cinemática direta permite determinar a posição e a orientação do \textit{end effector} de um manipulador robótico a partir das suas variáveis articulares. No presente projeto, esta abordagem é utilizada para calcular a posição do centro de preensão do myCobot~280, expressa pelas coordenadas $(x,y,z)$, com base nos ângulos das seis juntas $(q_1,\ldots,q_6)$ e nas dimensões dos elos. Para uma configuração articular conhecida, o modelo cinemático direto permite determinar uma única pose do manipulador.

O myCobot~280 é um manipulador constituído por seis juntas rotativas. Cada junta permite uma rotação em torno de um determinado eixo, enquanto os elos estabelecem as distâncias entre os eixos de rotação. Na configuração de referência adotada, correspondente a todas as juntas a $0^\circ$, o braço encontra-se na vertical e a garra na horizontal, orientada segundo o sentido positivo do eixo $y$.

Nesta configuração, a junta J1 roda em torno do eixo $z$, as juntas J2, J3 e J4 rodam em torno de eixos paralelos a $x$, a junta J5 roda em torno de um eixo paralelo a $z$ e a junta J6 permite a rotação da garra em torno do seu próprio eixo longitudinal. Estas direções correspondem à configuração de referência, podendo os eixos das juntas alterar a sua orientação relativamente à base durante o movimento.

Seguindo a abordagem geométrica estudada nas aulas, o movimento do manipulador é decomposto em duas componentes: o posicionamento num plano vertical e a rotação desse plano em torno do eixo vertical da base.

Cada junta desempenha uma função específica:

\begin{itemize}
    \item A \textbf{junta 1} permite rodar o braço em torno do eixo vertical, definindo a orientação do plano no qual ocorre o movimento;

    \item As \textbf{juntas 2 e 3} permitem aplicar o modelo geométrico de um braço planar com dois graus de liberdade, determinando a posição da junta J4, correspondente ao pulso;

    \item A \textbf{junta 4} determina a orientação do elo $L_4$ no plano do braço, contribuindo para o posicionamento da junta J5 e do centro de preensão;

    \item A \textbf{junta 5} permite alterar a orientação da garra relativamente ao plano do braço. Inicialmente, considera-se $q_5=0^\circ$, sendo posteriormente analisado o efeito da sua rotação;

    \item A \textbf{junta 6} permite rodar a garra em torno do seu eixo longitudinal. Considerando o centro de preensão alinhado com este eixo, a sua rotação altera a orientação da garra, sem modificar a posição do centro de preensão.
\end{itemize}

A Figura~\ref{fig:3d} apresenta uma representação tridimensional do manipulador, identificando os elos, os eixos de rotação das juntas e as grandezas geométricas utilizadas no cálculo da posição do centro de preensão. Esta posição é descrita através de uma distância $d$ no plano do braço, de uma altura $h$ e de um desvio lateral associado à geometria do manipulador.

% Vista 3D do braço
\begin{figure}[H]
    \centering
    \resizebox{0.5\textwidth}{!}{
        \begin{tikzpicture}[x={(0.10cm,-0.11cm)}, y={(0.26cm,0.04cm)}, z={(0.00cm,0.25cm)}, >=stealth,
                    link/.style={line width=2.4pt, line cap=round, join=round},
                    sombra/.style={line width=2.4pt, line cap=round, gray!30},
                    eixo/.style={->, thin, red!70!black},
                    cota3/.style={dashed, thick, blue!60!black},
                    junta/.style={circle, fill=blue!50, draw=blue!40!black, inner sep=2.2pt}]
  % Coordenadas
    \coordinate (J1) at (0.0,0.0,0.0);
    \coordinate (J2) at (0.0,0.0,13.2);
    \coordinate (J3) at (-1.9,3.3,23.5);
    \coordinate (O) at (3.7,6.5,23.5);
    \coordinate (J4) at (-1.0,14.7,21.9);
    \coordinate (J5) at (-4.7,21.0,21.9);
    \coordinate (J6) at (-4.7,21.0,17.3);
    \coordinate (G) at (-4.7,21.0,6.7);

  % Eixos
    \fill[gray!8] (-12,-6,0) -- (14,-6,0) -- (14,30,0) -- (-12,30,0) -- cycle;
    \draw[->] (0,0,0) -- (14,0,0) node[below] {$x$};
    \draw[->] (0,0,0) -- (0,30,0) node[right] {$y$};
    \draw[->] (0,0,0) -- (0,0,34) node[above] {$z$};

  % Projeção
   \draw[sombra] (0.0,0.0,0) -- (0.0,0.0,0) -- (-1.9,3.3,0) -- (3.7,6.5,0) -- (-1.0,14.7,0) -- (-4.7,21.0,0) -- (-4.7,21.0,0) -- (-4.7,21.0,0);
  
  % Base
    \fill[gray!35] plot[domain=0:360, samples=48, variable=\t] ({4*cos(\t)},{4*sin(\t)},0);
    \fill[gray!25] plot[domain=0:360, samples=48, variable=\t] ({4*cos(\t)},{4*sin(\t)},3);
    \draw[gray!60!black] plot[domain=0:360, samples=48, variable=\t] ({4*cos(\t)},{4*sin(\t)},3);
    \draw[gray!60!black] (1.4,3.8,0) -- (1.4,3.8,3);
    \draw[gray!60!black] (-1.4,-3.8,0) -- (-1.4,-3.8,3);
  
  % Braços
    \draw[link] (J1) -- node[left] {$L_1$} (J2);
    \draw[link] (J2) -- node[left] {$L_2$} (J3);
    \draw[link] (O) -- node[above] {$L_3$} (J4);
    \draw[link] (J4) -- node[above left] {$L_4$} (J5);
    \draw[link] (J5) -- node[right] {$L_6$} (J6);
    \coordinate (palma) at (-4.7,21.0,9.3);
    \draw[link] (J6) -- node[right] {$L_7$} (palma);
 
  % Garra
    \draw[line width=2pt, line cap=round] (-5.7,22.8,9.3) -- (-3.7,19.3,9.3);
    \draw[line width=2pt, line cap=round] (-5.6,22.6,9.3) -- (-5.6,22.6,5.5);
    \draw[line width=2pt, line cap=round] (-3.8,19.5,9.3) -- (-3.8,19.5,5.5);
  
  % Desvio lateral L5 (ao longo do eixo de J3)
    \draw[line width=2.4pt, line cap=round, orange!85!black] (J3) -- (O);
    \draw[thin, gray] (J3) -- (-1.9,3.3,29.0) (O) -- (3.7,6.5,29.0);
    \draw[<->, thin] (-1.9,3.3,28.7) -- node[above=7pt, align=center, font=\small] {$L_5$\\desvio lateral} (3.7,6.5,28.7);
  
  % Cotas
    \coordinate (S) at (5.6,3.2,0.0);
    \coordinate (F) at (-4.7,21.0,0.0);
    \draw[cota3] (0,0,0) -- node[pos=0.6, below=1pt] {$L_5$} (S);
    \draw[cota3] (S) -- node[below=2pt] {$d$} (F);
    \draw[cota3] (F) -- node[pos=0.5, right=2pt] {$h$} (G);
    \draw[blue!60!black, thin] (4.9,4.4,0.0) -- (3.8,3.7,0.0) -- (4.5,2.6,0.0);
    \draw[blue!60!black, thin] (-4.7,21.0,1.3) -- (-4.0,19.9,1.3) -- (-4.0,19.9,0.0);
  
  % Eixos de rotação das juntas
    \draw[eixo] (0.0,0.0,-2.8) -- (0.0,0.0,3.2);
    \draw[eixo] (-2.4,-1.4,13.2) -- (2.8,1.6,13.2);
    \draw[eixo] (-6.2,0.8,23.5) -- (3.2,6.2,23.5);
    \draw[eixo] (-3.4,13.3,21.9) -- (1.8,16.3,21.9);
    \draw[eixo] (-3.3,18.6,21.9) -- (-6.3,23.8,21.9);
    \draw[eixo] (-4.7,21.0,20.1) -- (-4.7,21.0,14.1);
  
  % Juntas
    \node[junta] at (J1) {};
    \node[junta] at (J2) {};
    \node[junta] at (J3) {};
    \node[junta] at (J4) {};
    \node[junta] at (J5) {};
    \node[junta] at (J6) {};
    \node[right=6pt] at (J1) {J1 (base)};
    \node[left=6pt] at (J2) {J2 (ombro)};
    \node[below left=2pt and 4pt] at (J3) {J3 (cotovelo)};
    \node[below=7pt] at (J4) {J4 (pulso)};
    \node[right=7pt] at (J5) {J5};
    \node[right=6pt] at (J6) {J6};
    \fill (G) circle (1.5pt);
    \node[right=12pt, font=\small] at (G) {centro de preensão};
    \draw[eixo] ([shift={(-4.6cm,-0.2cm)}]current bounding box.north east) -- ++(0.7cm,0) node[right, black] {\small eixo de rotação da junta};

  \end{tikzpicture}
    }
    \caption{Representação tridimensional do manipulador myCobot~280.}
    \label{fig:3d}
\end{figure}

O cálculo da posição é realizado em duas etapas. Inicialmente, determina-se a posição no plano lateral, através das coordenadas $(d,h)$, considerando os ângulos $q_2$, $q_3$ e $q_4$, bem como a influência de $q_5$ na projeção da garra, conforme apresentado na Secção~\ref{sec:plano_lateral}.

Posteriormente, a rotação da junta J1 permite transformar essa posição para o referencial da base, obtendo as coordenadas tridimensionais $(x,y,z)$, como descrito na Secção~\ref{sec:plano_horizontal}.

Os parâmetros geométricos utilizados nesta formulação foram obtidos a partir do modelo do manipulador no CoppeliaSim e encontram-se apresentados na Tabela~\ref{tab:elos}.

\begin{table}[H]
    \centering
    \caption{Comprimentos dos elos, medidos no cenário do CoppeliaSim.}
    \label{tab:elos}
    \begin{tabular}{llc}
        \hline
        Símbolo & Ligação & Comprimento [m] \\
        \hline
        $L_1$ & base $\rightarrow$ J2 (altura do ombro) & 0,1316 \\
        $L_2$ & J2 $\rightarrow$ J3 & 0,1104 \\
        $L_3$ & J3 $\rightarrow$ J4 & 0,0960 \\
        $L_4$ & J4 $\rightarrow$ J5 & 0,0731 \\
        $L_5$ & desvio lateral J3 $\rightarrow$ J4 & 0,0646 \\
        $L_6$ & J5 $\rightarrow$ J6 & 0,0456 \\
        $L_7$ & J6 $\rightarrow$ centro de preensão & 0,1057 \\
        \hline
    \end{tabular}
\end{table}

Os comprimentos apresentados correspondem às medidas utilizadas no cenário de simulação, podendo apresentar pequenas diferenças relativamente aos valores fornecidos pelo fabricante.

O parâmetro $L_1$ representa a altura da segunda junta relativamente à base, enquanto $L_2$, $L_3$, $L_4$ e $L_6$ correspondem às distâncias entre juntas consecutivas. O parâmetro $L_5$ representa o desvio lateral da estrutura e $L_7$ corresponde à distância entre a última junta e o centro de preensão.

% Secções do capítulo

\section{Plano lateral}
\label{sec:plano_lateral}

A análise do plano lateral permite estudar o movimento do braço numa vista de lado, utilizando duas coordenadas: $d$, que representa a distância frontal, e $h$, que representa a altura do centro de preensão relativamente à base.

Nesta análise, considera-se $q_1=0^\circ$ como referência, uma vez que a rotação da base não altera as coordenadas $d$ e $h$. Quando J1 roda, o plano lateral acompanha essa rotação, sendo a sua influência considerada posteriormente no plano horizontal.

Como ilustrado na Figura~\ref{fig:lateral}, as juntas J2, J3 e J4 possuem eixos de rotação paralelos e perpendiculares ao plano de movimento. Assim, a posição do braço pode ser calculada através das projeções dos seus elos, utilizando uma abordagem geométrica planar.

\begin{figure}[H]
    \centering
    \resizebox{0.6\textwidth}{!}{%Vista lateral do robô
% (L4 desenhado inclinado para se ver o ângulo theta_C; a garra
%  é perpendicular a L4, por isso alpha = theta_C - 90 graus)
\begin{tikzpicture}[scale=0.3,
                    % Estilos
                        link/.style={line width=2.2pt, line cap=round},
                        junta/.style={circle, draw, fill=white, line width=1pt, inner sep=2.2pt},
                        cota/.style={gray!70!black, thin},
                        ang/.style={->, thin},
                        qang/.style={->, thin, red!75!black},
                        qcota/.style={red!60, thin, dashed}]

  % Coordenadas das juntas
    \coordinate (J1) at (0,0);
    \coordinate (J2) at (0,13.2);
    \coordinate (J3) at (3.8,23.5);
    \coordinate (J4) at (13.2,21.9);
    \coordinate (J5) at (20.06,24.40);
    \coordinate (J6) at (25.26,10.12);

  % Eixos
    \draw[->, thin] (-4,0) -- (31,0) node[below left] {$d$};
    \draw[->, thin] (J1) -- (0,29) node[left] {$h$};

  % Base
    \fill[gray!25] (-3,-0.8) rectangle (3,0);
    \draw[thick] (-5,0) -- (30.5,0);

  % Ângulos absolutos (theta_A, theta_B, theta_C e alpha, a partir da horizontal)
    \draw[cota, dashed] (J2) -- ++(6,0);
    \draw[ang] ($(J2)+(4,0)$) arc[start angle=0, end angle=69.75, radius=4] node[pos=0.55, right=1pt] {$\theta_A$};
    \draw[cota, dashed] (J3) -- ++(9,0);
    \draw[ang] ($(J3)+(8,0)$) arc[start angle=0, end angle=-9.66, radius=8] node[pos=0.5, right=2pt] {$\theta_B$};
    \draw[cota, dashed] (J4) -- ++(7,0);
    \draw[ang] ($(J4)+(4.5,0)$) arc[start angle=0, end angle=20, radius=4.5] node[pos=0.5, right=1pt] {$\theta_C$};
    \draw[cota, dashed] (J5) -- ++(5,0);
    \draw[ang] ($(J5)+(3.5,0)$) arc[start angle=0, end angle=-70, radius=3.5] node[pos=0.45, right=2pt] {$\alpha$};

  % Ângulos das juntas (q2, q3, q4, relativos ao elo anterior)
    \draw[qcota] (J2) -- ++(0,7.5);
    \draw[qang] ($(J2)+(0,6.5)$) arc[start angle=90, end angle=69.75, radius=6.5] node[pos=0.5, left=3pt] {$q_2 < 0$};
    \draw[qcota] (J3) -- ++(69.75:5);
    \draw[qang] ($(J3)+(69.75:4)$) arc[start angle=69.75, end angle=-9.66, radius=4] node[pos=0.35, above right=0pt] {$q_3 < 0$};
    \draw[qcota] (J4) -- ++(-9.66:5);
    \draw[qang] ($(J4)+(-9.66:3)$) arc[start angle=-9.66, end angle=20, radius=3];
    \node[red!75!black] at ($(J4)+(3.6,-3)$) {$q_4 > 0$};

  % Braços
    \draw[link] (J1)  -- node[pos=0.55, right] {$L_1$} (J2);
    \draw[link] (J2) -- node[right] {$L_2$} (J3);
    \draw[link] (J3) -- node[below left=0pt and -2pt] {$L_3$} (J4);
    \draw[link] (J4) -- node[above=1pt] {$L_4$} (J5);
    \draw[link] (J5) -- node[right=2pt] {$D = L_6 + L_7$} (J6);

  % Garra (perpendicular a L4)
    \begin{scope}[shift={(J6)}, rotate=20]
      \draw[line width=1.4pt] (-1.1,1.5) -- (-1.1,-0.5) (1.2,1.5) -- (1.2,-0.5);
      \draw[line width=1.4pt] (-1.1,1.5) -- (1.2,1.5);
    \end{scope}

  % Juntas
    \node[junta] at (J2) {};  \node[left=5pt]  at (J2) {J2 (ombro)};
    \node[junta] at (J3) {};  \node[above left=3pt and 0pt] at (J3) {J3 (cotovelo)};
    \node[junta] at (J4) {};  \node[below=5pt] at (J4) {J4 (pulso)};
    \node[junta] at (J5) {};  \node[above=4pt] at (J5) {J5};
    \fill (J6) circle (6pt);   \node[below right=4pt and 2pt] at (J6) {centro de preensão $(d, h)$};

  % Guias
    \draw[cota, dashed] (J6) -- (25.26,0);
    \draw[cota, dashed] (J6) -- (0,10.12);

\end{tikzpicture}}
    \caption{Vista lateral do manipulador no seu plano vertical.}
    \label{fig:lateral}
\end{figure}

\subsection{Ângulos absolutos}

Os ângulos articulares $q_i$ representam a rotação de cada junta relativamente à sua configuração de referência. No entanto, para calcular as projeções dos elos, é necessário conhecer a orientação de cada um relativamente à horizontal.

Na configuração de referência, quando $q_2=0^\circ$, o elo $L_2$ encontra-se na vertical, formando um ângulo de $90^\circ$ com a horizontal. Por esse motivo, a sua orientação absoluta é obtida somando $90^\circ$ ao ângulo da junta J2.

Os elos seguintes dependem também das rotações das juntas anteriores. Assim, obtêm-se:

\begin{equation}
    \theta_A=q_2+\frac{\pi}{2}, \qquad
    \theta_B=\theta_A+q_3, \qquad
    \theta_C=\theta_B+q_4, \qquad
    \alpha=\theta_C-\frac{\pi}{2}
    \label{eq:angulos}
\end{equation}

Os ângulos $\theta_A$, $\theta_B$ e $\theta_C$ representam as orientações absolutas dos elos $L_2$, $L_3$ e $L_4$, respetivamente.

Como o eixo longitudinal da garra é perpendicular ao elo $L_4$, a sua orientação no plano lateral, representada por $\alpha$, é obtida subtraindo $90^\circ$ a $\theta_C$.

Deste modo, quando $\alpha=-90^\circ$, a garra encontra-se orientada verticalmente para baixo.

\subsection{Posição de cada junta}

Conhecidas as orientações dos elos, é possível calcular a posição de cada junta através das suas projeções.

Para um elo de comprimento $L$ e orientação $\theta$, a projeção horizontal é dada por $L\cos\theta$ e a vertical por $L\sin\theta$.

O cálculo é realizado de forma sucessiva, começando na base e acrescentando a contribuição de cada elo.

A junta J2, correspondente ao ombro, encontra-se a uma altura $L_1$. A partir desta posição, o elo $L_2$ permite determinar J3, o elo $L_3$ permite determinar J4 e o elo $L_4$ permite determinar J5.

Assim, considerando a base como origem do referencial, obtêm-se:

\begin{align}
    \text{J2 (ombro):}\quad
    &(d_2,h_2)=(0,L_1)
    \label{eq:j2}\\
    \text{J3 (cotovelo):}\quad
    &(d_3,h_3)=
    (d_2+L_2\cos\theta_A,\;
    h_2+L_2\sin\theta_A)
    \label{eq:j3}\\
    \text{J4 (pulso):}\quad
    &(d_4,h_4)=
    (d_3+L_3\cos\theta_B,\;
    h_3+L_3\sin\theta_B)
    \label{eq:j4}\\
    \text{J5:}\quad
    &(d_5,h_5)=
    (d_4+L_4\cos\theta_C,\;
    h_4+L_4\sin\theta_C)
    \label{eq:j5}
\end{align}

Estas expressões mostram que a posição de cada junta depende da posição anterior e da orientação do elo que as une.

Em particular, as juntas J2, J3 e J4 formam um sistema equivalente a um manipulador planar de dois graus de liberdade, estudado nas aulas \cite{bicho_2dof}.

Considerando $L_A=L_2$, $L_B=L_3$, $\theta_1=\theta_A$ e $\theta_2=q_3$, a posição de J4 relativamente a J2 pode ser escrita como:

\begin{equation}
    \begin{aligned}
        x_2 &= L_A\cos\theta_1+
        L_B\cos(\theta_1+\theta_2)\\
        y_2 &= L_A\sin\theta_1+
        L_B\sin(\theta_1+\theta_2)
    \end{aligned}
    \label{eq:2dof}
\end{equation}

Esta equivalência permite utilizar posteriormente a resolução geométrica de um manipulador de dois graus de liberdade na cinemática inversa, para determinar os ângulos necessários à posição pretendida de J4.

\subsection{Posição do centro de preensão}

As expressões anteriores permitem calcular a posição das juntas, mas o objetivo da cinemática direta é determinar a posição do centro de preensão da garra.

Para isso, é necessário acrescentar a extensão existente entre J5 e o centro de preensão.

Como os segmentos $L_6$ e $L_7$ estão alinhados com o eixo longitudinal da garra, esta distância é dada por:

\begin{equation}
    D=L_6+L_7
    \label{eq:distancia_garra}
\end{equation}

Considerando inicialmente $q_5=0^\circ$, a extensão da garra permanece no plano lateral, com orientação definida pelo ângulo $\alpha$.

Assim, a posição do centro de preensão é obtida somando as projeções dos elos $L_2$, $L_3$, $L_4$ e da extensão $D$.

As coordenadas resultantes são:

\begin{align}
    d &= L_2\cos\theta_A+
    L_3\cos\theta_B+
    L_4\cos\theta_C+
    D\cos\alpha
    \label{eq:d}\\
    h &= L_1+
    L_2\sin\theta_A+
    L_3\sin\theta_B+
    L_4\sin\theta_C+
    D\sin\alpha
    \label{eq:h}
\end{align}

Na primeira equação, as projeções horizontais dos elos determinam a distância frontal $d$.

Na segunda, somam-se as projeções verticais à altura inicial $L_1$, obtendo-se a altura $h$ do centro de preensão.

A Figura~\ref{fig:soma_projecoes} representa esta decomposição geométrica.

\begin{figure}[htbp]
    \centering
    \resizebox{0.6\textwidth}{!}{\begin{tikzpicture}[scale=0.3, >=stealth,
    elo/.style={line width=2.2pt, line cap=round},
    junta/.style={circle, draw, fill=white, line width=1pt, inner sep=2pt},
    guia/.style={gray!60, thin, densely dotted},
    seg/.style={line width=2.4pt},
    every node/.style={font=\footnotesize}]
  \colorlet{cL1}{gray!70!black}
  \colorlet{cL2}{blue!70!black}
  \colorlet{cL3}{orange!85!black}
  \colorlet{cL4}{green!50!black}
  \colorlet{cD}{violet!80!black}

  % Coordenadas
  \coordinate (B)  at (0,0);
  \coordinate (J2) at (0,13.2);
  \coordinate (J3) at (4.7,23.2);
  \coordinate (J4) at (14.1,21.5);
  \coordinate (J5) at (21.0,24.0);
  \coordinate (G)  at (26.2,9.8);

  % Eixos
  \draw[->, thin] (-2,0) -- (36,0) node[below left] {$d$};
  \draw[->, thin] (B) -- (0,28) node[left] {$h$};
  \fill[gray!25] (-2.5,-0.8) rectangle (2.5,0);

  % Guias
  \foreach \P in {J3,J4,J5,G} \draw[guia] (\P) -- (\P |- 0,-2.2);
  \draw[guia] (J2) -- (29,13.2);
  \draw[guia] (J3) -- (30,23.2);
  \draw[guia] (J4) -- (31,21.5);
  \draw[guia] (J5) -- (32,24.0);
  \draw[guia] (G)  -- (33,9.8);

  % Braços
  \draw[elo, cL1] (B)  -- (J2);
  \draw[elo, cL2] (J2) -- node[left]  {$L_2$} (J3);
  \draw[elo, cL3] (J3) -- node[pos=0.55, above] {$L_3$} (J4);
  \draw[elo, cL4] (J4) -- node[above left] {$L_4$} (J5);
  \draw[elo, cD]  (J5) -- node[right] {$D$} (G);
  \foreach \P in {J2,J3,J4,J5} \node[junta] at (\P) {};
  \fill (G) circle (6pt);
  \node[left=4pt]  at (J2) {J2};
  \node[above=3pt] at (J3) {J3};
  \node[below=4pt] at (J4) {J4};
  \node[above=3pt] at (J5) {J5};
  \node[below left=2pt] at (G) {garra};

  % Ângulos
  \draw[guia] (J2) -- ++(3.5,0);  \draw[->, thin, cL2] ($(J2)+(2.5,0)$) arc[start angle=0, end angle=65, radius=2.5] node[pos=0.4, right] {$\theta_A$};
  \draw[guia] (J3) -- ++(3.5,0);  \draw[->, thin, cL3] ($(J3)+(2.8,0)$) arc[start angle=0, end angle=-10, radius=2.8] node[pos=1, below right=-1pt] {$\theta_B$};
  \draw[guia] (J4) -- ++(3.5,0);  \draw[->, thin, cL4] ($(J4)+(2.8,0)$) arc[start angle=0, end angle=20, radius=2.8] node[pos=0.6, right] {$\theta_C$};
  \draw[guia] (J5) -- ++(3.5,0);  \draw[->, thin, cD]  ($(J5)+(2.5,0)$) arc[start angle=0, end angle=-70, radius=2.5] node[pos=0.6, right] {$\alpha$};

  % Soma de d
  \draw[seg, cL2] (0,-2.2)    -- (4.7,-2.2);
  \draw[seg, cL3] (4.7,-2.2)  -- (14.1,-2.2);
  \draw[seg, cL4] (14.1,-2.2) -- (21.0,-2.2);
  \draw[seg, cD]  (21.0,-2.2) -- (26.2,-2.2);
  \node[cL2, below=3pt] at (2.35,-2.2)  {$L_2\cos\theta_A$};
  \node[cL3, above=3pt] at (9.4,-2.2)   {$L_3\cos\theta_B$};
  \node[cL4, below=3pt] at (17.55,-2.2) {$L_4\cos\theta_C$};
  \node[cD,  above=3pt] at (23.6,-2.2)  {$D\cos\alpha$};
  \draw[|<->|, thin] (0,-5.2) -- node[fill=white, inner sep=1pt] {$d$} (26.2,-5.2);

  % Soma de h
  \draw[seg, cL1, ->] (29,0)    -- (29,13.2);
  \draw[seg, cL2, ->] (30,13.2) -- (30,23.2);
  \draw[seg, cL3, ->] (31,23.2) -- (31,21.5);
  \draw[seg, cL4, ->] (32,21.5) -- (32,24.0);
  \draw[seg, cD,  ->] (33,24.0) -- (33,9.8);
  \draw[guia] (29,13.2) -- (30,13.2) (30,23.2) -- (31,23.2) (31,21.5) -- (32,21.5) (32,24.0) -- (33,24.0);
  \node[cL1, left=1pt]  at (29,6.6)   {$L_1$};
  \node[cL2, right=1pt] at (30,17.5)  {$L_2\sin\theta_A$};
  \node[cL3, right=1pt] at (33.2,22.4) {$L_3\sin\theta_B < 0$};
  \node[cL4, right=1pt] at (32,25.6)  {$L_4\sin\theta_C$};
  \node[cD,  right=1pt] at (33,15.0)  {$D\sin\alpha < 0$};
  \draw[|<->|, thin] (36,0) -- node[fill=white, inner sep=1pt] {$h$} (36,9.8);
  \draw[guia] (33,9.8) -- (36,9.8);
  
\end{tikzpicture}}
    \caption{Decomposição dos elos e soma das projeções no plano lateral.}
    \label{fig:soma_projecoes}
\end{figure}

\subsection{Juntas 5 e 6}

Até este ponto, considerou-se $q_5=0^\circ$, mantendo a extensão da garra no plano lateral. No entanto, quando J5 roda, parte dessa extensão passa a encontrar-se fora do plano.

Nesse caso, a componente de $D$ que permanece no plano lateral é dada por $D\cos q_5$, enquanto a componente perpendicular corresponde a $D\sin q_5$.

Por isso, nas expressões anteriores, a extensão $D$ é substituída pela sua projeção $D\cos q_5$. Obtêm-se, assim:

\begin{align}
    d &= L_2\cos\theta_A+
    L_3\cos\theta_B+
    L_4\cos\theta_C+
    D\cos q_5\cos\alpha
    \label{eq:d_q5}\\
    h &= L_1+
    L_2\sin\theta_A+
    L_3\sin\theta_B+
    L_4\sin\theta_C+
    D\cos q_5\sin\alpha
    \label{eq:h_q5}\\
    l &= L_5-D\sin q_5
    \label{eq:q5}
\end{align}

O parâmetro $l$ representa o desvio lateral total do centro de preensão, considerando o desvio estrutural $L_5$ e a contribuição da rotação de J5.

Na implementação da cinemática inversa desenvolvida neste projeto, considera-se $q_5=0^\circ$. Consequentemente, a garra permanece no plano lateral e o desvio lateral reduz-se a $l=L_5$.

Por sua vez, J6 permite rodar os dedos da garra em torno do eixo longitudinal do \textit{end effector}. Como o centro de preensão se encontra alinhado com este eixo, esta rotação altera apenas a orientação dos dedos, sem modificar a posição do centro de preensão.

Assim, J6 não intervém no cálculo das coordenadas da posição.



\section{Plano horizontal}
\label{sec:plano_horizontal}

Após calcular a distância frontal $d$ e a altura $h$ no plano lateral, é necessário considerar a rotação da junta J1 para determinar a posição tridimensional do centro de preensão.

A junta J1 roda o braço em torno do eixo vertical $z$, alterando a direção do seu movimento no plano horizontal $(x,y)$, como ilustrado na Figura~\ref{fig:topo}.

\begin{figure}[htbp]
    \centering
    \resizebox{0.6\textwidth}{!}{%Vista de cima do robô
\begin{tikzpicture}[scale=0.3,
                    % Estilos
                    link/.style={line width=2pt, line cap=round},
                    cota/.style={gray!70!black, thin}]

  % Coordenadas das juntas
    \coordinate (J1) at (0,0);
    \coordinate (J6) at (-4.7,21);     % centro de preensão (x, y)
    \coordinate (O5) at (5.6,3.2);

  % Base
    \draw[->] (-10,0) -- (12,0) node[below left] {$x$};
    \draw[->] (0,-3) -- (0,25) node[left] {$y$};
    \fill[gray!25] (0,0) circle (2.5);  \draw (0,0) circle (2.5);
    \node[below left] at (-2.2,-0.2) {J1};
  
  % Desvio lateral L5 e plano do braço
    \draw[link] (J1) -- node[pos=0.55, below right=-2pt] {$L_5$} (O5);
    \draw[link] (O5) -- node[pos=0.4, right=2pt] {$d$} (J6);
    \draw[cota, dashed] (O5) -- ($(O5)!1.1!(J6)$) node[above, black] {plano do braço};
    \draw[cota, dashed] (O5) -- (30:14);
    \fill (J6) circle (7pt); \node[left=4pt] at (J6) {$(x, y)$};
  
  % Distância
    \draw[cota, dashed] (J1) -- node[pos=0.8, left=1pt, black] {$r$} (J6);
  
  % Ângulos
    \draw[->, thick] (5,0) arc[start angle=0, end angle=30, radius=5] node[pos=0, below right=2pt and -3pt] {$q_1 > 0$};
    \draw[->, thin] (8.5,0) arc[start angle=0, end angle=102.6, radius=8.5] node[pos=0.27, right=2pt] {$\operatorname{atan2}(y, x)$};
    \draw[->, thin] (30:12) arc[start angle=30, end angle=102.6, radius=12] node[pos=0.3, above right=0pt] {$\operatorname{atan2}(d, L_5)$};
    \draw[thin] ($(O5)+(120:1.3)$) -- ($(O5)+(120:1.3)+(210:1.3)$) -- ($(O5)+(210:1.3)$);

  % Convenção de sinal (visto de cima)
  %  \draw[->, thick, red!70!black] (9,23) ++(-30:1.6) arc[start angle=-30, end angle=250, radius=1.6];
  %  \node[red!70!black] at (9,23) {$+$};
  %  \node[right=2pt, align=left, font=\small] at (10.8,23) {ângulo positivo:\\sentido anti-horário};

\end{tikzpicture}}
    \caption{Vista superior do manipulador no plano horizontal.}
    \label{fig:topo}
\end{figure}

\subsection{Desvio lateral}

O manipulador robótico apresenta um desvio lateral $L_5$ de aproximadamente $6{,}5$ cm, devido à disposição dos seus elos.

Isto significa que o plano de movimento do braço não passa exatamente pelo eixo de rotação da base, encontrando-se deslocado lateralmente, como representado na Figura~\ref{fig:topo}.

Na formulação desenvolvida, este desvio é representado por $l$, que inclui também o efeito da rotação da junta J5, conforme definido na Secção~\ref{sec:plano_lateral}.

Assim, a posição horizontal do centro de preensão depende de duas distâncias perpendiculares:
\begin{itemize}
    \item $d$: distância frontal, calculada no plano lateral;
    \item $l$: desvio lateral relativamente ao eixo da base.
\end{itemize}

\subsection{Rotação da base}

Na configuração de referência, quando $q_1=0^\circ$, a distância frontal $d$ encontra-se orientada segundo o eixo positivo $y$, enquanto o desvio lateral $l$ está orientado segundo o eixo positivo $x$.

Por isso, nesta configuração, as coordenadas do centro de preensão são:

\begin{equation}
    x=l,\qquad y=d,\qquad z=h
\end{equation}

Quando J1 roda de um ângulo $q_1$, as direções de $d$ e $l$ também rodam em torno do eixo vertical $z$.

Para determinar as novas coordenadas, projetam-se estas duas distâncias nos eixos $x$ e $y$.

A componente lateral $l$ contribui com $l\cos q_1$ para $x$ e $l\sin q_1$ para $y$.

Por sua vez, a componente frontal $d$, perpendicular a $l$, contribui com $-d\sin q_1$ para $x$ e $d\cos q_1$ para $y$.

Somando estas contribuições, obtêm-se as coordenadas tridimensionais do centro de preensão:

\begin{equation}
    \boxed{
    \begin{aligned}
        x &= l\cos q_1-d\sin q_1\\
        y &= l\sin q_1+d\cos q_1\\
        z &= h
    \end{aligned}
    }
    \label{eq:xyz}
\end{equation}

A coordenada $z$ mantém-se igual a $h$, porque a rotação de J1 ocorre em torno do eixo vertical e não altera a altura da garra.

Na ausência de desvio lateral ($l=0$), as expressões simplificam-se para:

\begin{equation}
    x=-d\sin q_1,\qquad
    y=d\cos q_1,\qquad
    z=h
\end{equation}

Neste caso, o plano de movimento do braço passa pelo eixo de rotação da base.

\subsection{Distância ao eixo da base}

Conhecidas as coordenadas $x$ e $y$, é possível determinar a distância horizontal $r$ entre o centro de preensão e o eixo da base.

Como as distâncias $d$ e $l$ são perpendiculares, formam um triângulo retângulo, pelo que se aplica o teorema de Pitágoras:

\begin{equation}
    r=\sqrt{x^2+y^2}=\sqrt{l^2+d^2}
    \label{eq:r}
\end{equation}

Esta relação permite determinar a distância frontal a partir de uma posição pretendida, sendo utilizada posteriormente na cinemática inversa.

Além da distância, é necessário conhecer a direção do alvo no plano horizontal.

O ângulo formado pela posição $(x,y)$ relativamente ao eixo $x$ é dado por $\operatorname{atan2}(y,x)$.

Este ângulo corresponde à soma da rotação da base $q_1$ com o ângulo formado pelas componentes $l$ e $d$:

\begin{equation}
    \operatorname{atan2}(y,x)
    =
    q_1+\operatorname{atan2}(d,l)
    \label{eq:angulo_horizontal}
\end{equation}

Esta expressão será utilizada na cinemática inversa para determinar o ângulo $q_1$ necessário para orientar o manipulador na direção de um alvo.


\section{Resumo e verificação}
\label{sec:dk_resumo}

Combinando os resultados obtidos nos planos lateral e horizontal, é possível determinar a posição tridimensional do centro de preensão do myCobot~280 em função dos ângulos articulares e dos parâmetros geométricos do manipulador.

Considerando $D=L_6+L_7$ e os ângulos absolutos definidos na Equação~\eqref{eq:angulos}, a formulação da cinemática direta é dada por:

\begin{equation}
    \boxed{
    \begin{aligned}
        d &= L_2\cos\theta_A
        + L_3\cos\theta_B
        + L_4\cos\theta_C
        + D\cos q_5\cos\alpha\\
        h &= L_1
        + L_2\sin\theta_A
        + L_3\sin\theta_B
        + L_4\sin\theta_C
        + D\cos q_5\sin\alpha\\
        l &= L_5-D\sin q_5\\[4pt]
        x &= l\cos q_1-d\sin q_1\\
        y &= l\sin q_1+d\cos q_1\\
        z &= h
    \end{aligned}
    }
    \label{eq:dk_final}
\end{equation}

Estas expressões permitem calcular a posição do centro de preensão relativamente ao referencial da base, a partir dos ângulos das juntas e dos comprimentos dos elos.

Adicionalmente, a direção longitudinal da garra pode ser determinada através das componentes $-\sin q_5$, $\cos q_5\cos\alpha$ e $\cos q_5\sin\alpha$, expressas no referencial associado ao plano do braço. Estas componentes são posteriormente transformadas para o referencial da base, considerando a rotação $q_1$.

\subsection{Verificação dos resultados}

Para verificar o comportamento do modelo matemático, a cinemática direta foi avaliada em duas configurações distintas do manipulador.

\subsubsection{Configuração de referência}

A primeira configuração corresponde à posição inicial do manipulador, designada por \textit{pose zero}, na qual todas as juntas apresentam um ângulo de $0^\circ$.

De acordo com a Equação~\eqref{eq:angulos}, nesta configuração obtêm-se:

\begin{equation}
    \theta_A=\theta_B=\theta_C=90^\circ,
    \qquad \alpha=0^\circ
\end{equation}

Assim, o braço encontra-se orientado verticalmente, enquanto a garra permanece na horizontal, apontando no sentido positivo do eixo $y$.

As expressões da cinemática direta simplificam-se para:

\begin{equation}
    d=L_6+L_7, \qquad
    h=L_1+L_2+L_3+L_4
    \qquad \Rightarrow \qquad
    (x,y,z)=(L_5,d,h)
    \label{eq:pose_zero}
\end{equation}

Substituindo os valores dos comprimentos apresentados na Tabela~\ref{tab:elos}, obtém-se:

\begin{equation}
    (x,y,z)=(0{,}0646;\ 0{,}1513;\ 0{,}4111)
    \ \text{m}
\end{equation}

A direção longitudinal da garra é representada pelo vetor $(0,1,0)$, estando o resultado de acordo com a geometria definida para a configuração de referência.

\subsubsection{Configuração de trabalho}

A segunda configuração corresponde a uma posição de trabalho, na qual a garra se encontra orientada verticalmente para baixo.

Os ângulos articulares considerados são:

\begin{equation}
    \mathbf{q}=
    (18{,}84^\circ;\
    -12{,}74^\circ;\
    -94{,}21^\circ;\
    16{,}95^\circ;\
    0^\circ;\
    0^\circ)
    \label{eq:q_trabalho}
\end{equation}

Aplicando as relações definidas para os ângulos absolutos, obtêm-se:

\begin{equation}
    \theta_A=77{,}26^\circ,\qquad
    \theta_B=-16{,}95^\circ,\qquad
    \theta_C=0^\circ,\qquad
    \alpha=-90^\circ
\end{equation}

Substituindo estes valores nas equações do plano lateral, obtêm-se as coordenadas:

\begin{equation}
    d=0{,}1893\ \text{m},\qquad
    h=0{,}0600\ \text{m}
\end{equation}

Considerando a rotação da base e o desvio lateral do manipulador, a posição tridimensional do centro de preensão é:

\begin{equation}
    (x,y,z)\approx
    (0{,}0000;\ 0{,}2000;\ 0{,}0600)
    \ \text{m}
\end{equation}

Deste modo, o centro de preensão encontra-se aproximadamente a $20$ cm de distância horizontal da base, na direção positiva do eixo $y$, e a uma altura de $6$ cm.

A Figura~\ref{fig:poses_fk} apresenta as duas configurações consideradas na verificação da cinemática direta.

\begin{figure}[H]
    \centering
    \resizebox{0.75\textwidth}{!}{
        % Pose zero e configuração de trabalho (plano do braço); unidades em cm
\begin{tikzpicture}[scale=0.26, >=stealth,
    elo/.style={line width=2pt, line cap=round},
    junta/.style={circle, draw, fill=white, line width=0.9pt, inner sep=1.8pt},
    cota/.style={gray!70!black, thin},
    ang/.style={->, thin},
    every node/.style={font=\footnotesize}]

  \draw[->, thin] (-3,0) -- (22,0) node[below left] {$d$};
  \draw[->, thin] (0,0) -- (0,46) node[left] {$h$};
  \fill[gray!25] (-2.5,-0.8) rectangle (2.5,0.0);
  \draw[elo, blue!70!black] (0.0,0.0) -- (0.0,13.2) -- (0.0,24.2) -- (0.0,33.8) -- (0.0,41.1) -- (15.1,41.1);
  \node[junta] at (0.0,13.2) {};
  \node[junta] at (0.0,24.2) {};
  \node[junta] at (0.0,33.8) {};
  \node[junta] at (0.0,41.1) {};
  \fill (15.1,41.1) circle (5pt);
  \node[right=4pt, align=left] at (13.1,37.1) {garra a apontar para $+y$\\$(x, y, z) = (0{,}0646;\ 0{,}1513;\ 0{,}4111)$ m};
  \node[font=\small\bfseries] at (8,50) {Pose zero};
  \draw[->, thin] (35,0) -- (62,0) node[below left] {$d$};
  \draw[->, thin] (38,0) -- (38,46) node[left] {$h$};
  \fill[gray!25] (35.5,-0.8) rectangle (40.5,0.0);
  \draw[elo, orange!85!black] (38.0,0.0) -- (38.0,13.2) -- (40.4,23.9) -- (49.6,21.1) -- (56.9,21.1) -- (56.9,6.0);
  \node[junta] at (38.0,13.2) {};
  \node[junta] at (40.4,23.9) {};
  \node[junta] at (49.6,21.1) {};
  \node[junta] at (56.9,21.1) {};
  \fill (56.9,6.0) circle (5pt);
  \node[below=4pt, align=center] at (56.9,6.0) {garra a apontar para baixo\\$(x, y, z) = (0;\ 0{,}2000;\ 0{,}0600)$ m};
  \node[font=\small\bfseries] at (50,50) {Configuração de trabalho};
\end{tikzpicture}
    }
    \caption{Configuração de referência e configuração de trabalho utilizadas na verificação da cinemática direta.}
    \label{fig:poses_fk}
\end{figure}

\subsubsection{Validação no CoppeliaSim}

A validação do modelo cinemático é realizada através do modo de diagnóstico implementado em MATLAB, comparando a posição do centro de preensão calculada pela cinemática direta com a posição obtida diretamente no CoppeliaSim.

O erro de posicionamento é determinado através da distância euclidiana entre as duas posições, sendo expresso em milímetros:

\begin{equation}
    e_{\mathrm{modelo}} =
    1000
    \left\|
    \mathbf{p}_{\mathrm{DK}} -
    \mathbf{p}_{\mathrm{sim}}
    \right\|_2
    \label{eq:erro_modelo}
\end{equation}

onde $\mathbf{p}_{\mathrm{DK}}$ representa a posição calculada pelo modelo de cinemática direta e $\mathbf{p}_{\mathrm{sim}}$ corresponde à posição do centro de preensão obtida no CoppeliaSim, ambas expressas em metros e relativamente ao referencial da base.

No modo de diagnóstico, o programa obtém os ângulos articulares do manipulador, calcula a posição através da função de cinemática direta e compara o resultado com a posição fornecida pelo simulador.

Nos resultados reportados para as configurações testadas, as diferenças observadas situaram-se entre $1{,}2$ mm e $4{,}0$ mm, indicando uma boa aproximação entre o modelo geométrico desenvolvido e o comportamento do manipulador em simulação.

Adicionalmente, o programa permite calcular o erro entre a posição pretendida e a posição efetivamente alcançada, possibilitando avaliar o desempenho do sistema de posicionamento.

Foi ainda considerada a consistência entre a cinemática direta e a cinemática inversa, verificando se a aplicação sucessiva dos dois modelos permite recuperar a posição inicialmente pretendida.

\clearpage

